\section{Methods in full}\label{app:methods}

Every procedure the experiments compare against is specified here, so the paper
stands without the repository. Each subsection names the file it was run from and
that file's SHA-256 at the time of the runs reported.

\subsection{Intrinsic self-assessment: metrics and evidence scopes}
\label{app:methods:intrinsic}

Source: \texttt{src/probes.py}, sha256
\texttt{0c99cac7240ec9153c5417756b2cf6d8d35e0c647ae6253649250162a5d1f7cd}.

Every intrinsic reading is a second call to the same model that produced the step,
given the agent's own pre-action context plus the step's reasoning and proposed
action. Three metrics were run; all three are reported as \emph{uncertainty} $U$,
with $1$ meaning maximally uncertain, so higher always means more likely
incorrect.

\paragraph{$\ptrue$.} The model is asked a closed question --- ``Is the proposed
action above the correct and appropriate next action for this task? Answer with a
single word: Yes or No.'' --- and the answer is read from the \emph{first token's}
top-logprobs rather than from generated text. Probability mass is summed over
case and leading-space variants of Yes and of No, and $U = P(\mathrm{No}) /
(P(\mathrm{Yes}) + P(\mathrm{No}))$. A record with no yes/no mass in the head is
marked unparsed and excluded; it is never imputed. This is the only metric that
needs logprob access, which is why a frontier model that refuses logprobs can be
the agent in this framework but not the estimator.

\paragraph{Separate verbalized confidence.} The model is asked, in a separate
call, how confident it is that the step is correct, and answers in words. The
reply is parsed to a number in $[0,1]$ and $U = 1 - \mathrm{confidence}$. This is
the ``grounded verbalized channel'' the survey literature hopes for; on our
measurements it is the inverted one for capable models.

\paragraph{Post-hoc numeric confidence.} The same question on an explicit
$0$--$100$ scale (``On a scale of 0-100, how confident are you that the agent's
response above \ldots''), parsed to an integer and mapped to
$U = 1 - n/100$. Unparseable replies are excluded on the same rule.

\paragraph{Probe scope.} Every headline result uses a single whole-step
probe: one call about the step as a unit, reasoning together with the chosen
action. Narrower variants (reasoning-only, action-only, and their mean) exist
in the release and underlie the recipe-stability analysis of
\S\ref{sec:instrument}; no headline number uses them. Probe
seeds are derived per (step, probe, stage) so no two probe calls on a step share
an RNG stream.

\subsection{The supervised process critic}\label{app:methods:critic}

Source: \texttt{analysis/rev\_critic.py}, sha256
\texttt{b9331ffd0001ad405712c1bbf659095f9a4288d8d355137bc2fa31404ae4bd74}.

A per-step binary classifier trained on cheap trajectory features, with no
additional model calls, evaluated per capable external cell on the primary
construct.

\paragraph{Features (leak-free set).} Step index and $\log(1+\mathrm{index})$;
episode length so far; action string length in characters and in words; a one-hot
action type over the sixteen ALFWorld/HotpotQA verbs
(\texttt{go, take, open, close, put, toggle, heat, cool, clean, examine,
inventory, look, search, lookup, finish, none}); plan length before the step, the
change in plan length, and a plan-missing indicator; an environment indicator; and
a one-hot over agent identity. The \emph{hybrid} variant adds history-only
summaries of the external estimator's score over steps strictly before $t$
(mean, max, min, count, last value), which cannot leak $U_t$.

\paragraph{What is excluded, and why it must be.} Five features are withheld from
the leak-free set: \texttt{in\_admissible}, \texttt{empty\_result},
\texttt{loop\_flag}, \texttt{repeat\_dist}, \texttt{repeat\_count}. Under the
primary construct these are not predictors of the label --- they are the Tier-A
rules that \emph{define} its positive class, so a critic given them recomputes the
label rather than predicting it. Including them moves the comparison by more than
$0.12$ balanced accuracy and produces an apparent critic win; the leak-free set is
the honest comparison and is the one reported.

\paragraph{Models and hyperparameters, fixed in advance.} Logistic regression
(\texttt{C}$=1.0$, \texttt{max\_iter}$=2000$, lbfgs solver) on standardized
features, and gradient-boosted trees (100 estimators, max depth 3, learning rate
$0.1$). No hyperparameter search was run at any point.

\paragraph{Budget protocol.} Episodes are sorted by identifier and split in half;
both fold directions are run and averaged. Within the fitting half, steps are
taken \emph{in episode order} --- not sampled --- until the budget is reached,
giving budgets of exactly $200$ and $1{,}000$ labelled steps. The estimator's own
cut is fitted on the identical steps, so the two sides of the comparison see the
same labels. Scoring is out-of-sample on the other half, and the $1{,}000$-step
row covers only the cells whose fitting half contains that many labelled steps.

\subsection{The unsupervised mixture cut}\label{app:methods:mixture}

Scores are mapped to the logit scale and a two-component univariate Gaussian
mixture is fitted by EM: $k$-means initialisation on perturbed quartiles, three
restarts keeping the best log-likelihood, at most 200 iterations, convergence at
$10^{-6}$. The cut is the equal-posterior boundary between the two components,
found on a 4001-point grid between the fitted means (the closed-form quadratic
root can fall outside the data range), then mapped back through the logistic.

\paragraph{Abstention rule.} The arm declines to place a cut, rather than placing
a bad one, in four cases: EM failure; a minority component weight below $0.05$; a
separation between component means below $0.25$ on the logit scale; or no sign
change of the posterior difference between the means. Abstentions are counted as
losses wherever a pass rate is reported --- never dropped.

The procedure uses no labels at all, only the target's own score distribution.

\subsection{The fixed-percentile and error-rate-indexed rules}
\label{app:methods:simplerules}

Sources: \texttt{the released cut-transfer module}, sha256
\texttt{3021afcb3a716bd4a4985a77dc888daeb14334cf6f1ca907650877e5ade36a1e};
\texttt{the released indexed-cut module}, sha256
\texttt{9463a268b096b146f8429d7a111b39b311815a7398cafc1cd69356e543f9e91b}.

Both rules transfer from the same reference pool the operating-point map uses: all
other targets of that (estimator, environment), self-cells excluded.

\paragraph{Raw threshold transfer.} Pool the reference cells' steps, fit the
balanced-accuracy-optimal cut on the pool, and apply that \emph{score value}
directly to the new target. This is the baseline the level-shift result rules out.

\paragraph{Fixed percentile.} Take the same pooled cut, record what percentile it
occupies within the pooled scores, and place the new target's cut at that
percentile of its own unlabelled scores. Transfers a rank rather than a level, and
needs no labels on the target.

\paragraph{Error-rate-indexed rule.} Fit a linear map $g$ from a reference cell's
mean score to its error rate, over the same pool. For a new target, predict
$\hat{p} = g(\bar{u})$, clipped to $[0.02, 0.98]$, and cut at the
$100(1-\hat{p})$-th percentile of its own scores. This embeds the identity
``optimal percentile $= 100(1-p)$'', which the measurements contradict: realised
optimal percentiles deviate from it by up to $26$ points, so even an oracle handed
the true error rate places the cut suboptimally.

\paragraph{Common evaluation.} All rules place a cut without reading any label of
the target being scored, and all are evaluated by balanced accuracy against the
same labels on the same steps. Realised flag rates match requested percentiles to
within $0.001$; no tie-lumping occurs.
